\documentclass[11pt]{article}

\usepackage[utf8]{inputenc}
\usepackage[T1]{fontenc}
\usepackage{hyperref}
\usepackage{url}
\usepackage{booktabs}
\usepackage{amsfonts}
\usepackage{amsmath}
\usepackage{microtype}
\usepackage{graphicx}
\usepackage{natbib}
\usepackage{xcolor}
\usepackage[margin=1in]{geometry}

\newcommand{\kaos}{\textsc{Kaos}}

\title{Benchmark Gifts: Measurable Workload Properties That Decide\\Whether Agent Mechanisms Transfer}

\author{
  Danilo Canivel \\
  Independent / \kaos{} project \\
  \texttt{d.canivel@gmail.com}
}

\begin{document}
\maketitle

\begin{abstract}
Mechanisms for self-improving LLM agents---skill libraries, episodic memories, error-correction graphs, harness adaptation layers---are published at a growing rate, almost always evaluated on closed benchmarks. Practitioners who adopt them face an unanswered question: will the mechanism survive contact with a real workload? We report evidence from a longitudinal evaluation program in \kaos{}, an open-source local-first agent harness, in which thirteen published mechanism candidates were evaluated over six months and none was shipped without surviving a pre-registered, hash-locked, falsifiable probe. Two full probes, two rejection-by-measurement studies, and a four-axis measurement of the harness's organic workload converge on a simple account: benchmarks quietly supply four properties---closed action vocabularies, dense outcome signals, lexically anchored task text, and per-task expert references---that we call \emph{benchmark gifts}, and a mechanism transfers only if the axes its premise consumes are actually present in the target workload. In our organic workload, two gifts are measurably absent (89.0\% of agent trajectories normalize to a single repeated action label; five plasticity tables hold zero organically produced rows), one is absent at the grain mechanisms need (exact expert-pair coverage 6.8\%), and one is present (85.4\% of task texts carry hard lexical anchors)---and the ledger matches axis by axis: a graph-diff mechanism scored 1.000 against a 0.188 baseline on constructed workload yet failed the absent axis its premise consumes, while retrieval mechanisms consuming the present axis produced the program's only shipped gains. We contribute the axes, the measurement protocol, a probe lifecycle case study in which a pre-registered gate passed vacuously and was caught by instrument audit, and a practical recommendation: audit workload shape before building.
\end{abstract}

\section{Introduction}

In July 2026 we evaluated a mechanism that scored a perfect 1.000 on its evaluation workload, beat the incumbent baseline by 81 percentage points, passed every one of its four pre-registered kill gates---and was not shipped. The mechanism, a contrastive trajectory-diff localizer distilled from a published error-correction paper \citep{emg2026}, had been probed under the falsifiable-evaluation discipline that governs \kaos{} development: kill gates written and cryptographically hash-locked before any feature code existed, a mandatory falsification self-test, and a binding verdict with no retune-and-rerun \citep{canivel2026plasticity}. The gate that should have protected against exactly this outcome---a requirement that real, organically produced agent trajectories exhibit the node-reuse structure the mechanism's premise consumes---\emph{passed}, at a median of 10.0 against a floor of 1.30. A post-run instrument audit showed why: 97 of 109 organic agents run one tool in a loop, so their trajectories normalize to a \emph{single repeated label}, satisfying a one-sided threshold vacuously. Among the twelve agents whose trajectories are diverse enough for a diff to operate on at all, median reuse was 1.000 and zero of twelve reached the floor. The mechanism's premise is false on the workload it was intended for; the verdict on file remains ACCEPT, the disposition is do-not-ship, and a regression test now fails if the mechanism ever enters the shipped package (\S\ref{sec:gdl}).

This paper is about the pattern that episode belongs to. Over six months and thirteen published mechanism candidates, the \kaos{} evaluation program has repeatedly watched mechanisms die not on their own logic but on the \emph{shape of the workload}: an action-validation layer that could never be evaluated because the organic workload produced two qualifying incidents against a floor of two hundred \citep{lifeharness2026}; a retrieval-side consolidation family rejected because lexical-only retrieval cannot simultaneously satisfy token-faithfulness and query-bridging \citep{canivel2026plasticity}; and now a trajectory-graph mechanism whose closed-vocabulary premise dissolves against real tool-argument entropy. Meanwhile the mechanisms that \emph{did} survive measurement---BM25-based retrieval re-ranked by usage statistics---consume the one workload property our measurements show real tasks actually supply in abundance: hard lexical anchors.

We name the underlying variables \emph{benchmark gifts}: properties that evaluation environments provide implicitly and gratis, that mechanism designs consume as unstated premises, and that real workloads may or may not supply. We propose four, all cheaply measurable:

\begin{itemize}
  \item \textbf{Contributions.}
  \item \emph{The four-axis workload-shape audit} (\S\ref{sec:axes}): closed action vocabulary (M1), outcome-signal density (M2), task-text lexical anchoring (M3), and expert-reference availability (M4), each with a deterministic, embedding-free measurement procedure, applied to a real harness's organic database (Table~\ref{tab:axes}).
  \item \emph{A longitudinal, pre-registered evaluation ledger} (\S\ref{sec:ledger}): thirteen published mechanism candidates evaluated under one constraint set with explicit evidence tiers---to our knowledge the first such series run under hash-locked pre-registration.
  \item \emph{A probe lifecycle case study} (\S\ref{sec:gdl}) in which a pre-registered gate passed vacuously, demonstrating a failure mode of pre-registration itself (one-sided predicates) and the instrument-audit practice that catches it.
  \item \emph{A falsifiable synthesis} (\S\ref{sec:analysis}): mechanism families mapped to the axes they consume, with the prediction---testable by any practitioner in an afternoon---that a mechanism fails to transfer when its consumed axis is absent.
\end{itemize}

Everything quantitative in this paper traces to a committed artifact in the public \kaos{} repository via a provenance ledger (Appendix~\ref{app:provenance}); claim labels (W1--W13) refer to its rows.

\section{Related work}

\textbf{Do published agent mechanisms help in practice?} The question has begun to receive direct negative evidence. SWE-Skills-Bench found 39 of 49 deployed agent skills produced zero measurable gain on real software-engineering tasks, with three degrading performance by up to 10 points \citep{sweskills2026}. CTIM-Rover found cross-task episodic memory failed to beat its baseline in any tested configuration, with retrieved items actively harming exploration \citep{ctimrover2025}. HASP's own ablation showed unfiltered skill evolution costs 24 points \citep{hasp2026}. These studies establish \emph{that} transfer fails; our contribution is a measured account of \emph{which workload properties} decide it, plus a discipline that catches the failure before shipping rather than after.

\textbf{Harness engineering and the verification gap.} The 2026 discourse converged on the position that the harness, not the model, is the active frontier \citep{hashimoto2026journey, osmani2026harness, zhou2026externalization}, with the sharpest critique being the absence of verification that harness changes help \citep{bockeler2026harness}. Recent systems work proposes evolution control planes that select among memory, skill, and weight interventions \citep{areal2026}, but leaves the decision objective unoperationalized; \citet{metaharness2026} optimize harnesses end-to-end against benchmark objectives, inheriting the benchmark-gift problem this paper measures. Evaluation methodology work has begun to quantify variance in agentic evals \citep{multirun2026}; we add the orthogonal axis: even a variance-correct benchmark result transfers only conditional on workload shape.

\textbf{Pre-registration and falsification.} Our discipline is Popperian pre-commitment applied to software mechanisms \citep{popper1959logic}: kill gates are written and hash-locked before feature code exists, a falsification self-test must prove the harness \emph{can} reject the feature, and verdicts are binding. The companion paper describes the discipline and the plasticity architecture it governs \citep{canivel2026plasticity}; this paper treats the discipline as instrument and reports what six months of its output reveals about transfer. Classic self-improvement systems---Voyager's skill library \citep{wang2023voyager}, Reflexion's verbal reinforcement \citep{shinn2023reflexion}---are upstream of the mechanism families evaluated here.

\section{Setting and method}

\subsection{The harness and its constraint set}

\kaos{} is an open-source, local-first multi-agent harness whose entire runtime state resides in one SQLite file, with hard constraints relevant here: no embeddings, no parameter updates of any kind, no GPU requirement, no mandatory external services \citep{canivel2026plasticity}. The constraints matter methodologically: they force every adopted mechanism into a deterministic, auditable form, which is what makes the workload-shape measurements of \S\ref{sec:axes} exact rather than sampled.

\subsection{Evidence tiers}

We report three tiers and label every claim: \textbf{Tier A}, pre-registered hash-locked probes with binding verdicts (two run to completion, two rejection-by-measurement studies); \textbf{Tier B}, deterministic mechanical measurements without a lock (the workload-shape audit; there is no mechanism under test, hence nothing to falsify); \textbf{Tier C}, structured panel evaluations of external papers (used only for motivation and candidate selection, never as evidence for a quantitative claim). The thirteen-candidate ledger (Table~\ref{tab:ledger}) spans all three tiers, and we consider the tier separation itself part of the method: a frequent failure of survey-style critique is treating read-the-paper judgments as if they were measurements.

\subsection{The four axes}
\label{sec:axes}

Each axis names a property that benchmark environments supply implicitly, together with the deterministic procedure we use to measure its presence in an organic workload (a live \kaos{} database: 453 agents, 1{,}639 tool calls, 171 task texts, accumulated by ordinary development and ARC-competition operation; composition disclosed in \S\ref{sec:limits}).

\textbf{M1 --- Closed action vocabulary.} Text-game benchmarks admit exhaustive action normalization (e.g., \texttt{(take, apple, table)}); graph- and pattern-based mechanisms consume this by assuming repeated visits to identical normalized actions. \emph{Measure:} normalize every tool call to a label---tool name plus a value-class signature over its arguments under fixed string rules (paths bucketed by top directory and extension, volatile keys dropped, long strings bucketed; frozen in the GDL probe lock, W5)---and compute per-agent \emph{node reuse}: calls divided by distinct labels. \emph{Result (W1):} 89.0\% of agents with $\geq$8 calls are mono-label ($<$3 distinct labels); among the 12 label-diverse agents, median reuse is 1.000 (p75 = 1.120) and \textbf{0\%} reach the 1.30 floor the GDL probe pre-registered. Organic trajectories are either one action repeated or all-unique chains; the middle regime that trajectory-graph mechanisms need does not occur.

\textbf{M2 --- Outcome-signal density.} Benchmarks emit a reward every episode and often every step. \emph{Measure:} count organically produced rows in the harness's outcome-carrying tables, and the tool-call error rate. \emph{Result (W2):} the five plasticity tables (\texttt{episode\_signals}, \texttt{skill\_uses}, \texttt{failure\_fingerprints}, \texttt{critical\_steps}, \texttt{memory}) contain \textbf{zero} organically produced rows; 2 of 1{,}639 tool calls (0.12\%) errored; the only outcome truth is episode-grain agent status (295 completed / 97 failed / 24 killed). This corroborates at population scale what a 2026 probe found at incident scale: an action-validation mechanism's binding evaluation returned VOID because the organic workload had produced $n{=}2$ qualifying incidents against a pre-registered floor of 200 (W6).

\textbf{M3 --- Task-text lexical anchoring.} Benchmark task descriptions are templated, making task-similarity retrieval trivial for embeddings and hard for lexical methods---an inversion of the usual assumption. \emph{Measure:} the fraction of task texts containing at least one hard lexical anchor (file path, snake\_case/camelCase identifier, version number, error code, hex hash), by fixed regular expressions. \emph{Result (W3):} \textbf{85.4\%} (146/171) of organic task texts are anchored (paths in 144, identifiers in 127, versions in 124), with a median length of 656 tokens. Real operational tasks name their objects. This is the one gift the organic workload supplies \emph{more} generously than benchmarks do---and it is consumed by exactly the mechanism family that has survived measurement in this program (W9).

\textbf{M4 --- Expert-reference availability.} Correction-from-comparison mechanisms assume a per-task successful reference: EMG consumes one expert trajectory per training task from a curated dataset \citep{emg2026}. \emph{Measure:} for each failed agent with a recorded task text, whether a completed agent exists with an identical task text (exact), or with token-Jaccard $\geq 0.5$ (loose). \emph{Result (W4):} exact coverage \textbf{6.8\%} (5/73); loose coverage 68.5\% (50/73). Loose pairs are common---but M1 makes them undiffable, so the gift is absent at the grain the mechanism family needs.

\section{The ledger as a longitudinal study}
\label{sec:ledger}

Table~\ref{tab:ledger} summarizes the thirteen candidates evaluated since v0.7. Six were rejected, four parked with pre-registered re-entry conditions, one returned VOID (uninterpretable by lock design, a distinct and honest outcome), one was accepted-but-not-shipped (\S\ref{sec:gdl}), and the retrieval mechanisms that ship in \kaos{} today entered before this series and survived its measurement standards retroactively (W9). Zero mechanisms shipped on unverified promise. The right-hand column records the axis on which each evaluation turned, where one did.

\begin{table}[t]
\centering
\small
\begin{tabular}{@{}llll@{}}
\toprule
Candidate (source) & Tier & Outcome & Deciding axis \\
\midrule
SAGE graph memory & C & reject (constraint conflict) & --- \\
Synthesis-as-consolidation & A & \textbf{reject, measured} (W7) & M3$^{\dagger}$ \\
Extractive consolidation & A & \textbf{do-not-build, measured} (W7) & M3$^{\dagger}$ \\
AutoResearchClaw & C & orthogonal; 1 idea parked & --- \\
HASP skill programs & C & do-not-build (external priors) & M2 \\
Life-Harness action layer & A & \textbf{VOID \#1, organic $n{=}2 < 200$} (W6) & M2 \\
UserHarness & C & park & --- \\
Per-agent DB sharding & A/B & \textbf{reject by measurement} (W8) & --- \\
MATM transactive memory & C & reject (embeddings); TRUF parked & --- \\
ATG task graphs & C & park (6/10) & M1 (anticipated) \\
AReaL2.0 / ATDP & C & reject; 2 fragments distilled & M2 \\
EMG $\rightarrow$ GDL probe & A & \textbf{accept-per-lock, do-not-ship} (W5) & M1 \\
Retrieval re-ranking (ships) & A/B & \textbf{+10.0 / +13.3 pp} (W9) & M3 (present) \\
\bottomrule
\end{tabular}
\caption{Thirteen mechanism candidates evaluated 2026-01 through 2026-07 under one constraint set. Tier A = pre-registered hash-locked probe; B = deterministic measurement; C = panel evaluation (motivation only). Bold = measured outcome. $^{\dagger}$The synthesis rejections established the \emph{FTS vise}: under lexical-only retrieval a stored abstraction must be simultaneously token-faithful and query-bridgeable, which no tested variant satisfied (LLM arm 0.000 with the synthesized item present in top-$K$ $\approx$50\% of the time; extractive arm 0.000 despite 44/44 token faithfulness)---a structural boundary of the M3 axis rather than an absence of it.}
\label{tab:ledger}
\end{table}

Three regularities emerge. First, \emph{no mechanism died on its own internal logic}: every measured rejection traced to a workload property the mechanism's benchmark had supplied for free. Second, the two Tier-A probes that ran to completion both turned on their workload gates, not their capability gates---the capability claims (localization accuracy, validation lift) were either confirmed or never reached. Third, the single mechanism family with measured positive transfer consumes the single axis measured present.

\section{Case studies}

\subsection{GDL: a perfect score, a vacuous gate, and a do-not-ship ACCEPT}
\label{sec:gdl}

The Graph-Diff Localizer probe (W5) is, to our knowledge, the first published end-to-end record of a pre-registered agent-mechanism evaluation in which the verdict and the shipping decision correctly \emph{diverged}. The timeline, all commits public:

\begin{enumerate}
  \item \textbf{Lock first.} Kill gates frozen and sha256-pinned at a commit containing no feature code: G1, accuracy $\geq$ 0.70 absolute and $+$10pp over the shipped single-trajectory localizer, $n \geq 40$; G2, a wrong-pair lesion isolating the pairing's causal role; G3, \emph{organic} median node-reuse $\geq$ 1.30 (the reality gate); G4, lexical pairing precision $\geq$ 0.80. The workload was pre-registered as harness-generated (ground-truth divergence labels cannot exist organically), with the external-validity caveat written into the lock rather than discovered later.
  \item \textbf{Falsification self-test.} Substituting the feature arm by the baseline (FULL := B1) emitted the required kill; substituting the lesion likewise. A harness that cannot reject its feature is inadmissible; this one could.
  \item \textbf{Binding run.} FULL 1.000, B1 0.188, wrong-pair lesion 0.000 ($n{=}48$); pairing precision 1.000; organic median reuse 10.000. All gates passed; verdict \textbf{ACCEPT} as computed. The B1 figure is itself a finding: the shipped heuristic localizer returns no answer on all 24 silent wrong-branch episodes because it keys on error steps, mechanically confirming the blind spot the mechanism targeted.
  \item \textbf{Instrument audit.} The G3 median of 10.0 was implausibly high and was audited before the verdict was reported: 86\% of organic tool calls come from search-worker agents that call one tool in a loop with a volatile-scrubbed argument, normalizing to a single label. A mono-label trajectory satisfies a one-sided reuse floor \emph{vacuously}---and is exactly as undiffable as the all-unique chains the gate was written to catch. On the informative subpopulation (label diversity $\geq$ 3): median 1.000, \textbf{0/12} at the floor.
  \item \textbf{Resolution.} Under the no-retune clause the gate was not edited and the run was not repeated; the verdict stands as computed. The \emph{disposition} is do-not-ship: an ACCEPT whose only organic gate passed vacuously does not clear a shipping bar of ``no mechanism ships on hope.'' Shipping requires a successor lock whose reuse gate excludes mono-label trajectories---a gate that on today's data would kill. A regression test enforces that no graph-diff module exists in the shipped package while the disposition stands (W13).
\end{enumerate}

The methodological point: pre-registration constrains goalpost-moving but cannot guarantee a predicate measures what its author intended. The failure mode (one-sided thresholds satisfied by an unanticipated degenerate regime) is general, and the countermeasure---a post-run instrument audit that may overrule the \emph{disposition} while leaving the \emph{verdict} untouched---preserves both honesty and the no-retune property. We know of no equivalent practice in published agent-mechanism evaluations.

\subsection{VOID as data: the action-realization probe}

The Life-Harness probe (W6) pre-registered an organic-only workload with a floor of 200 action-class incidents and explicitly forbade synthetic substitution. The organic count was 2. The binding verdict was VOID---not REJECT: the mechanism is unevaluated, not refuted---and the lock's synthetic ban is what made the outcome informative: it converted ``we could not evaluate this'' from a silent methodology hole into a measured statement about signal scarcity (M2) on real workloads. Population-scale corroboration arrived five weeks later in the workload-shape audit: the outcome-signal tables that any such mechanism would learn from contain zero organic rows (W2).

\subsection{The positive control}

A transfer account that only explains failures explains nothing. The retrieval-re-ranking family---BM25 over verbatim text, re-ranked by Wilson-bound success statistics and recency \citep{wilson1927probable}---consumes M3 and only M3: it needs lexically anchored text on both sides of the match and episode-grain outcomes, the two things the organic workload measurably has. It produced this program's only shipped, measured gains ($+$10.0pp and $+$13.3pp top-1 over BM25 on small committed benchmarks, $n{=}10$/$n{=}15$, small-$n$ disclosed; W9), and the GDL probe's pairing gate independently measured its premise at 1.000 precision on anchored task text. The axis account predicts both columns of the ledger.

\section{Analysis: the audit, and what promotion gates miss}
\label{sec:analysis}

\textbf{The workload-shape audit.} The practical protocol this paper argues for costs an afternoon: before building a published mechanism into a production harness, measure the four axes on the target workload's existing telemetry (every measurement in \S\ref{sec:axes} is a read-only query plus fixed string rules), then read the mechanism's evaluation section and list the gifts its benchmark supplied. A mechanism whose consumed axis is absent will not transfer, no matter how large its headline delta---GDL's was 81 points. The audit is falsifiable: it predicts, for instance, that EMG-family mechanisms transfer to workloads with repeated task families and low argument entropy (regression-test suites, scheduled runbooks), and fail on heterogeneous coding/ops traffic; and that skill libraries whose retrieval keys are anchored verbatim artifacts will outperform ones keyed on synthesized summaries on any M3-positive workload (the FTS-vise result, W7, is this prediction's boundary case).

\textbf{Why promotion gates are not enough.} The prevailing verification proposals for self-evolving agents---shadow evaluation, replay, canary, rollback \citep{areal2026}---are \emph{promotion} gates: they ask whether a change helps on the traffic used to evaluate it. Every failure documented here would have sailed through them. GDL passes any promotion gate run on its constructed workload; the action-realization layer would have shipped with no evaluation at all (nothing in a promotion pipeline notices $n{=}2$); the synthesis mechanisms produced plausible artifacts that a canary would keep. The failures were caught by, respectively: a workload-shape gate plus instrument audit, a pre-registered sample-size floor with a synthetic ban, and kill gates written before the feature existed. The difference is Popperian: promotion gates accumulate confirmation; pre-registered probes purchase the right to ship by first proving the evaluation can refute.

\textbf{Instrument risk is part of the method's honest accounting.} Two of this program's most informative findings were defects in its own instruments: a 2026 self-audit found the verdict machinery could emit tautological agreement statistics and accept on empty gate sets \citep{canivel2026plasticity}, and the GDL gate passed vacuously (\S\ref{sec:gdl}). We report these not as embarrassments but as the discipline working: both were caught by audits the discipline mandates, both produced mechanical guards, and neither required weakening a verdict already on file.

\section{Limitations}
\label{sec:limits}

\textbf{One system, one workload.} All Tier A/B evidence comes from a single harness and its operator's organic database. The workload is disclosed but idiosyncratic: 86\% of tool calls come from ARC-competition search workers, which inflates the M1 mono-label rate; the label-diverse subpopulation behind the headline 0/12 figure has $n{=}12$. The axis \emph{measurements} are exact; their \emph{representativeness} is a single sample. The claim we defend is existential and predictive---these axes are measurable, they decided these thirteen outcomes, and the audit is cheap insurance---not a universal law of workloads.

\textbf{The axes are not proven exhaustive or independent.} Four axes suffice to explain this ledger; other workloads may fail mechanisms on axes we have not needed (e.g., non-stationarity). M1 and M4 interact (loose expert pairs exist but are undiffable), so the axes are diagnostic dimensions, not an orthogonal basis.

\textbf{Tier C is judgment, not measurement.} Six of thirteen ledger rows rest on structured panel evaluation of papers, not on probes. We use them only to characterize the candidate population and never cite a Tier-C number as evidence; still, candidate \emph{selection} is a human choice and a source of bias in which mechanisms got probes at all.

\textbf{M2's zeros are double-edged.} Empty outcome tables measure both organic signal scarcity and the fact that this development-era database predates routine use of the plasticity write paths. Both readings are true; the second is why delayed-outcome backfill remains a live (parked) candidate rather than a refuted one.

\textbf{Constructed-workload capability numbers do not claim external validity.} GDL's 1.000/0.188/0.000 figures were pre-registered as harness-generated and carry that caveat by design; we cite them as evidence about the \emph{baseline's blind spot} and the \emph{probe lifecycle}, not as expected field accuracy.

\section{Conclusion}

Across thirteen published mechanism candidates evaluated under pre-registered falsification, transfer was decided not by mechanism quality but by four measurable workload properties that benchmarks supply for free and real workloads may not. The practical consequence fits in one sentence: measure your workload's shape before you build---the queries are deterministic, the afternoon is cheap, and the alternative in our ledger was an 81-point benchmark win that would have shipped a mechanism whose premise occurs in zero of twelve real trajectories. Future work: replicating the axis measurements on other harnesses' telemetry (the protocol requires only tool-call logs and task texts); a successor GDL lock with a degeneracy-robust reuse gate; and probing the parked delayed-outcome backfill candidate, which targets the one axis (M2) that is absent today but is an artifact of write-path adoption rather than of the workload itself.

\section*{Acknowledgments and AI-use disclosure}

Consistent with the program it describes, this paper was drafted by an agent operating inside \kaos{}, against the provenance ledger of Appendix~\ref{app:provenance}, and revised through a multi-level review panel (grounding, rigor, adversarial) before human sign-off; panel records are committed alongside the source. All quantitative content derives from the committed artifacts cited in the ledger. Responsibility for all claims rests with the human author.

\bibliographystyle{plainnat}
\bibliography{references}

\appendix

\section{Provenance ledger}
\label{app:provenance}

Every quantitative claim in this paper carries a W-label resolving to a row of \texttt{paper2/PROVENANCE.md} in the \kaos{} repository (\url{https://github.com/canivel/kaos}), which maps the claim to a committed artifact (results file, lock file, verdict document, or commit hash) and an evidence tier. The workload-shape measurements (W1--W4) regenerate with \texttt{uv run python -m demo\_workload\_shape\_bench.bench}; the GDL probe (W5) with \texttt{uv run python -m demo\_graphdiff\_localizer\_bench.run}; the test suite state (W13) with \texttt{uv run python -m pytest tests/}.

\end{document}
